\documentclass[11pt]{article}
\usepackage{amsmath,amssymb,amsthm,graphicx,hyperref,booktabs,geometry}
\geometry{margin=1in}
\title{Black Holes as Almost-Perfect All:All Entanglement Graphs:\\Interior Collapse, Horizon Wiring, and the Micro-Hole Phase Transition}
\author{Draft v0.1 --- computational companion \\ \small\texttt{src/bh\_graph} + \texttt{scripts/generate\_figures.py} reproduce all figures}}
\date{\today}
\begin{document}
\maketitle
\begin{abstract}
We formalize the intuition that a black hole interior is an almost-perfect
all:all entanglement graph. Internal edges cost no exterior space; each of the
$k$ exterior legs costs $\sim l_p^2$ of horizon area, so $A(k)=k\,l_p^2$
independent of interior node count $N$. We show (1) the complete graph $K_N$
destroys interior distance (diameter~1, single-step spread) reproducing the
fast-scrambling hierarchy; (2) mass enters only via $k\propto M^2$, and the
$k\to0$ monogamy limit is a baby-universe pinch-off; (3) micro-holes are
pointlike until $k\,l_p^2 \ge 4\pi r_{\mathrm{point}}^2$, then pop a horizon ---
a phase transition with an LQG-style area gap. Simulations in
\texttt{bh\_graph} accompany each claim.
\end{abstract}

\section{Sec 1: All:all = no interior space}
Let interior wiring be $G_N$ on $N$ nodes. Compare $K_N$, chain $P_N$,
$\sqrt N\times\sqrt N$ grid, and random 3-regular graphs under deterministic SI
spread (cover time = seed eccentricity). Only $K_N$ has diameter~1 and cover
time~1 at every $N$ (Table~1, Figs.~1--2). This is the discrete shadow of
Sekino--Susskind fast scrambling and the ER=EPR reading that maximal
entanglement is zero-length connectivity.

\begin{table}[h]\centering
\begin{tabular}{lcccc}\toprule
family & diameter & $t_{\rm cover}$ & mean dist & gap \\\midrule
$K_N$ & 1 & 1 & 1 & $N$ \\
chain & $N{-}1$ & $\sim N/2$ & $\sim N/3$ & $\sim \pi^2/N^2$ \\
grid & $\sim 2\sqrt N$ & $\sim 2\sqrt N$ & $O(\sqrt N)$ & $O(1/N)$ \\
random-regular & $\sim\log N$ & $\sim\log N$ & $\sim\log N$ & $O(1)$ \\\bottomrule
\end{tabular}
\caption{Scrambling hierarchy; see \texttt{bh\_graph.scrambling}.}
\end{table}

\section{Sec 2: Horizon area counts exterior legs}
Write interior pairs $N(N{-}1)/2$ plus $k\ll N^2$ exterior legs. Postulate
\begin{equation}
A(k) = k\,l_p^2,\qquad R(k)=\sqrt{k\,l_p^2/4\pi},
\end{equation}
so $N$ drops out (Fig.~3). Schwarzschild consistency forces
$k(M)=4\pi(2M)^2/l_p^2\propto M^2$ (\texttt{k\_from\_mass\_schwarzschild}).
With toy monogamy frontier $e_{\rm int}+e_{\rm ext}\le 1$, $e_{\rm int}\to1$
implies $k\to0$: pinch-off into a closed graph / baby universe (Fig.~4).

\section{Sec 3: Micro-hole phase transition}
A point region of radius $r_{\rm point}$ embeds $k$ legs iff
$k\,l_p^2<4\pi r_{\rm point}^2$. Hence
\begin{equation}
R_{\rm obs}(k)=\begin{cases}r_{\rm point} & k\le k_{\rm crit}\\\sqrt{k\,l_p^2/4\pi} & k>k_{\rm crit}\end{cases},\quad k_{\rm crit}=4\pi r_{\rm point}^2/l_p^2,
\end{equation}
with LQG-style quantization $A\in\{0\}\cup[A_{\min},\infty)$ (Figs.~5--6).
Growth $k(N)=k_0+\alpha N$ stays pointlike then pops a horizon: micro-holes are
not scaled-down Schwarzschild holes.

\section*{Reproducibility}
\texttt{pip install -e .; python -m pytest tests/ -q; python scripts/generate\_figures.py; streamlit run app.py.}

\begin{thebibliography}{9}
\bibitem{ss08} Sekino \& Susskind, JHEP 2008 (fast scramblers).
\bibitem{vR10} Van Raamsdonk, Gen.\ Rel.\ Grav.\ 2010.
\bibitem{ms13} Maldacena \& Susskind, Fortsch.\ Phys.\ 2013 (ER=EPR).
\bibitem{bek73} Bekenstein, Phys.\ Rev.\ D 1973; Hawking 1975.
\bibitem{rs95} Rovelli \& Smolin, Nucl.\ Phys.\ B 1995.
\bibitem{isl19} Penington; Almheiri et al., JHEP 2019--2020 (islands).
\bibitem{hp07} Hayden \& Preskill, JHEP 2007.
\end{thebibliography}
\end{document}
